\paragraph{Differences among the four approaches in Section~1.2.}
D'Ascoli et al. directly sample a discrete lagged recurrence, execute it to
produce a sequence, and train a Transformer to recover the recurrence from that
sequence. Me\v{z}nar et al. learn an unconditional continuous prior over general
expression trees and search that latent space for equations fitting numerical
data. Bendinelli et al. retain an explicit random-tree prior for pretraining and
learn the conditional inverse map
$p_{\phi}(e\mid\mathcal D,\mathcal H)$, allowing hypotheses to steer which
equation is recovered. NSRwH is therefore a direct extension of the NeSymReS
family and contributes controlled recovery; D'Ascoli's construction supplies
the discrete temporal-generator role. ODEFormer retains random unary--binary
trees and prefix decoding, but interprets each tree as one component of a
continuous-time vector field and uses numerical integration to connect the
equation with observations. For this thesis, the four contributions compose
naturally: D'Ascoli supplies the lagged grammar and executable recurrence, HVAE
offers an optional learned prior over plausible recurrence trees, NSRwH supplies
a retriever that can be tested with different amounts and qualities of known
mechanism information, and ODEFormer supplies the continuous-time extension in
which evaluation must include trajectory reconstruction and generalization.

